\documentclass[10pt,a4paper]{amsart}
\usepackage[margin=19mm]{geometry}
\usepackage{amsmath,amssymb,mathtools}
\usepackage[scr=rsfs,cal=euler]{mathalfa}
\usepackage{xfrac}
\usepackage[unicode,hidelinks,bookmarks=false]{hyperref}
\newcommand{\ie}{i.e.\ }
\newcommand{\cf}{cf.\ }
\newcommand{\resp}{respectively\ }
\newcommand{\Subsection}{\subsection}
\providecommand{\nobreakdash}{\nobreak\hbox{-}}
\setlength{\emergencystretch}{2em}
\multlinegap0pt
\usepackage{bm}
\usepackage{bbm}
\usepackage{dsfont}
\usepackage{xspace}
\usepackage{cancel}
\usepackage{mathtools}
\usepackage{amsthm}
\makeatletter
\@ifundefined{thm}{\newtheorem{thm}{Theorem}}{}
\makeatother
\newcommand{\tIIAex}{\hat{\tilde{A}}}
\newcommand{\tIIAim}{\hat{A}}
\newcommand{\tIIaex}{\pmb{\hat{\tilde{a}}}}
\newcommand{\tIIaim}{\pmb{\hat{a}}}
\newcommand{\tIIbexf}{\hat{\tilde{b}}_1}
\newcommand{\tIIbex}{\pmb{\hat{\tilde{b}}}}
\newcommand{\tIIbimf}{\hat{b}_1}
\newcommand{\tIIbim}{\pmb{\hat{b}}}
\newcommand{\tIIcex}{\pmb{\hat{\tilde{c}}}}
\newcommand{\tIIcim}{\pmb{\hat{c}}}
\newcommand{\ukdv}{\eta}
\title[Traveling-wave solutions and structure-preserving numerical ]{Traveling-wave solutions and structure-preserving numerical methods for a hyperbolic approximation of the Korteweg-de Vries equation}
\author{Abhijit Biswas, David I. Ketcheson, Hendrik Ranocha, Jochen Sch\"utz}
\date{Source: arXiv:2412.17117v2 (CC BY 4.0)}
\begin{document}
\maketitle

\noindent\emph{Source and licence.} Traveling-wave solutions and structure-preserving numerical methods for a hyperbolic approximation of the Korteweg-de Vries equation. Version arXiv:2412.17117v2 (CC BY 4.0), distributed under the Creative Commons Attribution 4.0 International licence, as recorded in the arXiv licence metadata for this exact version and shown on the abstract page https://arxiv.org/abs/2412.17117v2. Full rights notice: licenses/arxiv-2412.17117-CC-BY-4.0.txt.\par\smallskip

\noindent\textit{Editorial scope.} Counted unit: the Thm whose source label is \texttt{Th:AP\_result\_type\_II} in the article (source0.tex lines 753--763, source label \texttt{Th:AP\_result\_type\_II}), delivered together with the article's own complete original proof of it (source0.tex lines 764--861, one \texttt{proof} environment of 98 lines). No mathematics is written, repaired or re-derived: statement and proof are the article's own text line for line. Required context retained uncounted: Eq:KdVH\_AP\_split (lines 527--545); Mthd:CK\_SA\_ARK (lines 572--586); hilbert (lines 640--642); Eq:well-prepared\_IC (lines 749--751). Every definition, display and label of those lines is retained unchanged. Omitted from this package and not redistributed: 1--526, 546--571, 587--639, 643--748, 752--752, 862--2213. Nothing inside a retained excerpt is omitted or altered: every retained body is delivered exactly as the article's lines are written, so no omission receipt of any kind accompanies this package. Citations: the retained excerpts carry no citation command at all, so no bibliography entry of the article is transcribed and no reference is redistributed. Reference adaptation: none was needed and none was made; every reference inside the delivered unit resolves inside it, so no unresolved reference or citation is produced. Printed slips kept literal: no printed word, symbol, hypothesis or number of the counted statement or of its proof was altered. Layout and class: the article's own class is \texttt{article} with packages including inputenc, amsthm, amsmath, amsfonts, hyperref, xcolor, float, geometry, graphicx, subcaption, caption, url, tikz; the wrapper is the lane's pinned \texttt{amsart} editorial wrapper at 10pt on A4, which restates the theorem-like environments the retained text needs and adds the article's own macro definitions that the retained text needs (\texttt{\textbackslash tIIAex}, \texttt{\textbackslash tIIAim}, \texttt{\textbackslash tIIaex}, \texttt{\textbackslash tIIaim}, \texttt{\textbackslash tIIbex}, \texttt{\textbackslash tIIbexf}, \texttt{\textbackslash tIIbim}, \texttt{\textbackslash tIIbimf}, \texttt{\textbackslash tIIcex}, \texttt{\textbackslash tIIcim}, \texttt{\textbackslash ukdv}), with extra wrapper packages (bm, bbm, dsfont, xspace, cancel, mathtools). The delivered numbering is the unit's own. Assets: no retained span contains a figure environment, a graphics-inclusion command, a PDF-page inclusion, a table or a TikZ picture; no image file is copied into this package, the package declares no asset dependency and no image file is redistributed with it. Licence: version arXiv:2412.17117v2 of this article is distributed under the Creative Commons Attribution 4.0 International licence, as recorded in the arXiv licence metadata for this exact version and shown on the abstract page https://arxiv.org/abs/2412.17117v2. A full rights notice is delivered at licenses/arxiv-2412.17117-CC-BY-4.0.txt. Characters: every retained excerpt is pure ASCII, so the delivered engine, XeLaTeX with its default Latin Modern text font, reports no missing character.
\begin{equation}\label{Eq:KdVH_AP_split}
    \underbrace{\begin{pmatrix}
        u \\
        v \\
        w
    \end{pmatrix}_{t}}_{= q_t}
    =
    \underbrace{\begin{pmatrix}
        -u u_x \\
        0 \\
        0
    \end{pmatrix}}_{= f(q)}
    +
    \underbrace{\begin{pmatrix}
        -w_x \\
        \frac{1}{\tau} (v_x - w) \\
        \frac{1}{\tau} (-u_x + v)
    \end{pmatrix}}_{= g(q)}.
\end{equation}

\begin{equation}\label{Mthd:CK_SA_ARK}
\begin{array}{c|ccc}
    0 & 0 \\
    \tIIcex & \tIIaex & \tIIAex\\
    \hline
    & \tIIbexf & \tIIbex^T
\end{array}
\qquad
\begin{array}{c|ccc}
    0 & 0 \\
    \tIIcim & \tIIaim & \tIIAim \\
    \hline
    & \tIIbimf & \tIIbim^T
\end{array} \;.
\end{equation}

\begin{align} \label{hilbert}
u^n = u^n_0 + \tau u^n_1 + \tau^2 u^n_2 + \cdots, \quad v^n  = v^n_0 + \tau v^n_1 + \tau^2 v^n_2 + \cdots, \quad w^n = w^n_0 + \tau w^n_1 + \tau^2 w^n_2 + \cdots \;,
\end{align}

\begin{equation}\label{Eq:well-prepared_IC}
    v^{0} = u^{0}_x +  \mathcal{O}(\tau) \quad \text{and} \quad w^{0} = u^{0}_{xx} +  \mathcal{O}(\tau).
\end{equation}

\begin{thm}\label{Th:AP_result_type_II}
    A globally stiffly accurate ImEx-RK method of type II, applied to the splitting \eqref{Eq:KdVH_AP_split} of the hyperbolic
    approximation of the KdV equation, together with the well-prepared initial data \eqref{Eq:well-prepared_IC}, is AP for all
    components $u$, $v$, and $w$. Furthermore, in the stiff limit $\tau \to 0$, the following error estimates apply to all
    components:
    \begin{align*}
        u^{n+1} - \ukdv(t_{n+1}) = \mathcal{O}(\Delta t^p), \ v^{n+1} - \ukdv_{x}(t_{n+1}) = \mathcal{O}(\Delta t^p),
        \textrm{and} \ w^{n+1} - \ukdv_{xx}(t_{n+1}) = \mathcal{O}(\Delta t^p)\;,
    \end{align*}
    where $p$ is the order of the ImEx-RK method.
\end{thm}

\begin{proof}
Let $q^n = [u^n, v^n, w^n]^T$ and $q^{n+1} = [u^{n+1}, v^{n+1}, w^{n+1}]^T$ denote the numerical approximations to the true solution at time $t_n$ and $t_{n+1} = t_n + \Delta t$, respectively.
Starting from the solution $q^n$ at time $t_n$, an ImEx-RK method of type II \eqref{Mthd:CK_SA_ARK} applied to the system
\begin{align*}
    q_t = f(q) + g(q)
\end{align*}
computes the approximate solution at time $t_{n+1}$ as
\begin{subequations}\label{Eq:ImEx_mthd}
\begin{align}
    q^{(1)} & = q^n \;, \\
    q^{(i)} & = q^n + \Delta t \left( \hat{\tilde{a}}_{i1} f(q^n) + \sum_{j = 2}^{i-1} \hat{\tilde{a}}_{ij} f(q^{(j)}) + \hat{a}_{i1} g(q^n) + \sum_{j = 2}^{i} \hat{a}_{ij} g(q^{(j)}) \right), \ i = 2,3, \ldots, s \;, \\
    q^{n+1} & = q^n + \Delta t \left( \tIIbexf f(q^n) + \sum_{j = 2}^{s} \hat{\tilde{b}}_{j} f(q^{(j)}) + \tIIbimf g(q^n) + \sum_{j = 2}^{s} \hat{b}_{j}  g(q^{(j)}) \right) \;.
\end{align}
\end{subequations}
Using the notation $\pmb{q} = [q^{(2)}, q^{(3)}, \dots, q^{(s)}]^T$, the identity matrix $I$ in $\mathbb{R}^{3 \times 3}$, and the vector of ones $\pmb{e}$ in $\mathbb{R}^{s-1}$, this can be written compactly as
\begin{subequations}
\begin{align}
    \pmb{q} &= \pmb{e} \otimes q^n + \Delta t \left( \tIIaex \otimes f(q^n) + (\tIIAex \otimes I) f(\pmb{q}) + \tIIaim \otimes g(q^n) + (\tIIAim \otimes I) g(\pmb{q}) \right) \;, \label{Eq:ImEx_stage} \\
    q^{n+1} &= q^n + \Delta t \left( \tIIbexf f(q^n) + (\tIIbex^T \otimes I) f(\pmb{q}) + \tIIbimf g(q^n) + (\tIIbim^T \otimes I) g(\pmb{q}) \right) \;. \label{Eq:ImEx_sol}
\end{align}
\end{subequations}
Let us denote by $\pmb{u} = [u^{(2)}, u^{(3)}, \dots, u^{(s)}]^T$, $\pmb{v} = [v^{(2)}, v^{(3)}, \dots, v^{(s)}]^T$, and $\pmb{w} = [w^{(2)}, w^{(3)}, \dots, w^{(s)}]^T$ the vectors of stage-solution components for the variables $u$, $v$, and $w$, respectively. For the KdVH system with the splitting \eqref{Eq:KdVH_AP_split}, the equation \eqref{Eq:ImEx_stage} in component form becomes
\begin{subequations}
\begin{align}
    \pmb{u} &= u^n \pmb{e} + \Delta t \left( (-u^n u_{x}^n)\tIIaex + \tIIAex (-\pmb{u} \pmb{u}_x)  -w_x^n \tIIaim + \tIIAim(-\pmb{w}_x) \right) \;, \label{Eq:stage_u} \\
    \pmb{v} &= v^n \pmb{e} + \frac{\Delta t}{\tau} \left( (v_x^n - w^n) \tIIaim + \tIIAim(\pmb{v}_x-\pmb{w})\right) \;, \label{Eq:stage_v} \\
    \pmb{w} &= w^n \pmb{e} + \frac{\Delta t}{\tau} \left( (-u_x^n + v^n) \tIIaim + \tIIAim(-\pmb{u}_x+\pmb{v})\right) \;. \label{Eq:stage_w}
\end{align}
\end{subequations}
Similarly, the final update of the solution can be written in components as
\begin{subequations}
\begin{align}
    u^{n+1} &= u^n  + \Delta t \left( \tIIbexf(-u^n u_{x}^n) +  \tIIbex^T  (-\pmb{u} \pmb{u}_x) + \tIIbimf(-w_x^n)  + \tIIbim^T  (-\pmb{w}_x)  \right) \;, \label{Eq:sol_u} \\
    v^{n+1} &= v^n  + \frac{\Delta t}{\tau} \left( \tIIbimf(v_x^n - w^n)  + \tIIbim^T  (\pmb{v}_x-\pmb{w}) \right) \;, \label{Eq:sol_v} \\
    w^{n+1} &= w^n + \frac{\Delta t}{\tau} \left( \tIIbimf(-u_x^n + v^n)  + \tIIbim^T  (-\pmb{u}_x+\pmb{v}) \right)  \;. \label{Eq:sol_w}
\end{align}
\end{subequations}
We assume there exist Hilbert expansions \eqref{hilbert} for $u^n$, $v^n$, and $w^n$
and similar expansions for the stage vectors $\pmb{u}$, $\pmb{v}$, and $\pmb{w}$ of the stage solution components:
\begin{align*}
\pmb{u} = \pmb{u}_0 + \tau \pmb{u}_1 + \tau^2 \pmb{u}_2 + \cdots, \quad \pmb{v}  = \pmb{v}_0 + \tau \pmb{v}_1 + \tau^2 \pmb{v}_2 + \cdots, \quad \pmb{w} = \pmb{w}_0 + \tau \pmb{w}_1 + \tau^2 \pmb{w}_2 + \cdots \;.
\end{align*}

We insert these expansions into the stage equations and compare the leading-order terms in the powers of $\tau$.
The leading-order terms in the expansions for the stage equations \eqref{Eq:stage_v} and \eqref{Eq:stage_w} yield:
\begin{align*}
    \mathcal{O}\left(\tau^{-1}\right) &\colon
    \begin{cases}
        \left((v_0^n)_x - w_0^n \right) \tIIaim + \tIIAim  \left((\pmb{v}_0)_x - \pmb{w}_0\right) = \pmb{0} \;, \\
        \left(-(u_0^n)_x + v_0^n \right) \tIIaim + \tIIAim  \left(-(\pmb{u}_0)_x + \pmb{v}_0\right) = \pmb{0} \;.
    \end{cases}
\end{align*}
The well-prepared initial data at time $t_n$ implies that $v_0^n = (u_0^n)_x$ and $w_0^n = (v_0^n)_x$.
Since $\tIIAim$ is invertible, it follows that
\begin{align}\label{Eq:Aux_var_rel}
   \pmb{v}_0 = (\pmb{u}_0)_x \quad \text{and} \quad \pmb{w}_0 = (\pmb{v}_0)_x \;.
\end{align}

The leading-order term in the expansion of \eqref{Eq:stage_u} gives
\begin{align*}
    \mathcal{O}\left(\tau^{0}\right) & \colon \pmb{u}_0 = u_0^n \pmb{e} + \Delta t \left( - u_0^n (u_{0}^n)_x \tIIaex - \tIIAex \pmb{u}_0 (\pmb{u}_0)_x -(w_0^n)_x \tIIaim - \tIIAim(\pmb{w}_0)_x \right)\;.
\end{align*}
Using $\pmb{v}_0 =  (\pmb{u}_0)_x$, $\pmb{w}_0 =  (\pmb{v}_0)_x$, and the well-preparedness of the initial data at time $t_n$
in the previous equation, we obtain
\begin{align}\label{Eq:Asymp_stage_u}
    \pmb{u}_0 & = u_0^n \pmb{e} + \Delta t \left(  - u_0^n (u_{0}^n)_x \tIIaex - \tIIAex \pmb{u}_0 (\pmb{u}_0)_x -(u_0^n)_{xxx} \tIIaim - \tIIAim(\pmb{u}_0)_{xxx}  \right)\;.
\end{align}

The solution update for the first equation in the leading-order term becomes:
\begin{align}\label{Eq:Asymp_sol_u}
    u_0^{n+1} &= u_0^n  + \Delta t \left(  - \tIIbexf u_0^n (u_{0}^n)_x  - \tIIbex^T  \pmb{u}_0 (\pmb{u}_0)_x -\tIIbimf(u_0^n)_{xxx}  - \tIIbim^T  (\pmb{u}_0)_{xxx} \right) \;.
\end{align}
In the limit as $\tau \to 0$, we have $u^n = u^n_0$ and $\pmb{u} = \pmb{u}_0$, so the numerical scheme becomes:
\begin{subequations}
    \begin{align*}
        \pmb{u}_0 & = u_0^n \pmb{e} + \Delta t \left(  - u_0^n (u_{0}^n)_x \tIIaex - \tIIAex \pmb{u}_0 (\pmb{u}_0)_x -(u_0^n)_{xxx} \tIIaim - \tIIAim(\pmb{u}_0)_{xxx}  \right)\;, \\
        u_0^{n+1} &= u_0^n  + \Delta t \left(  - \tIIbexf u_0^n (u_{0}^n)_x  - \tIIbex^T  \pmb{u}_0 (\pmb{u}_0)_x -\tIIbimf(u_0^n)_{xxx}  - \tIIbim^T  (\pmb{u}_0)_{xxx} \right)\;.
    \end{align*}
\end{subequations}
This is precisely the numerical scheme obtained from the time-stepping method \eqref{Mthd:CK_SA_ARK} when applied to the KdV
equation $\ukdv_t = - \ukdv \ukdv_x -\ukdv_{xxx} $, where the term $-\ukdv \ukdv_x$ is treated explicitly and the term $-\ukdv_{xxx}$ is treated implicitly.
It is important to note that, in order to preserve this result at the subsequent time step, the auxiliary components must
remain well-prepared. Therefore, we must project $v^{n+1}$ and $w^{n+1}$ onto the corresponding equilibrium manifolds. To achieve
this, we utilize the GSA property of the method. Since the method is assumed to satisfy the GSA property, it follows that:
\begin{subequations}
\begin{align*}
    v^{n+1} &= \tIIbim^T\tIIAim^{-1}\pmb{v} = v^{(s)} \;, \\
    w^{n+1} &= \tIIbim^T\tIIAim^{-1}\pmb{w} = w^{(s)} \;.
\end{align*}
\end{subequations}

In the limit as $\tau \to 0$, we have $v_0^{n+1} = v_0^{(s)}$ and $w_0^{n+1} = w_0^{(s)}$. Using equation
\eqref{Eq:Aux_var_rel}, we can express $v_0^{n+1} = (u_0^{(s)})_x$ and $w_0^{n+1} = (u_0^{(s)})_{xx}$.
By the GSA property of the scheme, it follows from equations \eqref{Eq:Asymp_stage_u} and \eqref{Eq:Asymp_sol_u}
that $u_0^{(s)} = u_0^{n+1}$, which leads to $v_0^{n+1} = (u_0^{n+1})_x$ and $w_0^{n+1} = (u_0^{n+1})_{xx}$.
This establishes the AP property, and the error estimates for all components follow trivially, thereby proving
the AA property of the method.
\end{proof}

\end{document}
